\section{Notation and diffusion conventions}
\label{app:notation}

\subsection{Symbols}

\begin{center}
\small
\begin{tabular}{ll@{\qquad}ll}
\toprule
symbol & meaning & symbol & meaning \\
\midrule
$\nsites$ & chain length $=$ ambient dimension & $\kernel,\ \kernelth$ & transition kernel \\
$\nseq$ & number of observed sequences & $\innov$ & innovation density \\
$\ngrid$ & number of quadrature grid points & $\corr$ & autoregressive coefficient \\
$\gridpt{1},\dots,\gridpt{\ngrid}$ & grid points & $\ivar$ & innovation variance \\
$\gridh$ & grid spacing & $\ncomp$ & mixture components \\
$\halfwidth$ & grid half-width, domain $[-\halfwidth,\halfwidth]$ & $\mu$ & initial law $\mu(a_1)$ \\
$t$ & diffusion time & $\like$ & unary likelihood factor \\
$\Dt = 1-e^{-2t}$ & noise variance & $\lmsg,\ \rmsg$ & forward, backward message \\
$\Pz,\ \Pt$ & clean, noised distribution & $\params$ & kernel parameters \\
$\score(x,t)$ & score $\nabla_x\log\Pt(x)$ & $\Qfun$ & EM auxiliary function \\
$\pmean{a}$ & posterior mean of $a$ given $x$ & $\Xis$ & expected transition statistic \\
$\gwn$ & standardised Gaussian white noise & $\Lik$ & marginal log-likelihood \\
$\tmin,\ \tmax$ & reverse-integration endpoints & $\lmmse$ & LMMSE estimator \\
\bottomrule
\end{tabular}
\end{center}

Two symbols are worth pinning down, because neighbouring literatures use them differently.
\textbf{$\kernel$ always means the transition kernel} and never a grid size; the grid size is
always $\ngrid$. The ambient dimension is written $\nsites$ because it \emph{is} the chain length
and every graphical-model statement is indexed by sites; the diffusion literature more often
writes $n$ for the sample count, which here is $\nseq$.

\subsection{Time orientation and the reverse SDE}

\begin{center}
\fbox{\parbox{0.93\textwidth}{
\textbf{Convention, used without variation.} $\gwn$ is standardised Gaussian white noise,
$\langle\gwn(t)\gwn(t')\rangle=\delta(t-t')$, so every $\sqrt{2}$ is written explicitly.

\emph{Forward}, $t$ running from $0$ to $\tmax$:
\[
\frac{\mathrm{d}x}{\mathrm{d}t} = -x + \sqrt{2}\,\gwn(t),
\qquad
x(t) = e^{-t}a + \sqrt{\Dt}\,z,
\qquad
\Dt = 1-e^{-2t},\quad z\sim\normal{0}{I}.
\]

\emph{Reverse}, in the elapsed-time variable $\tau = \tmax - t$ running from $0$ to
$\tmax-\tmin$, so that increasing $\tau$ means decreasing $t$:
\[
\frac{\mathrm{d}x}{\mathrm{d}\tau} = x + 2\,\score\!\left(x, \tmax-\tau\right) + \sqrt{2}\,\gwn(\tau).
\]

\emph{Probability-flow ODE}, same marginals, deterministic:
$\;\mathrm{d}x/\mathrm{d}\tau = x + \score(x,\tmax-\tau)$.

Initialisation $x(\tau{=}0)\sim\normal{0}{I}$; termination at $\tau = \tmax-\tmin$.
}}
\end{center}

The reverse drift follows from the time-reversal theorem for diffusions
\citep{anderson1982,haussmann1986}: for a forward SDE $\mathrm{d}x = f\,\mathrm{d}t +
g\,\mathrm{d}W$ the time-reversed process has drift $f - g^2\nabla\log p_t$ and the same
diffusion coefficient. With $f=-x$ and $g=\sqrt{2}$ the backward-in-$t$ drift is
$-x - 2\score$, and rewriting in $\tau$ flips the sign to give the box above. Because unit-variance
data keep $\Var(x_t)=1$, initialising at $\normal{0}{I}$ incurs a bias $e^{-2\tmax}$, which is
$2.5\times10^{-3}$ at $\tmax=3$. The implementation is
\codefile{src/reverse.py}, whose module docstring states the same convention; the sampler uses a
geometric time grid, denser at small $t$, because the reverse drift stiffens like
$1/\Dt$ there.

\section{The score--posterior identity}
\label{app:identity}

\paragraph{Differentiation under the integral.}
$\Pt(x) = \int\mathrm{d}a\,\Pz(a)\,g_{\Dt}(x-e^{-t}a)$ with
$g_{\Dt}(u) = (2\pi\Dt)^{-\nsites/2}e^{-\|u\|^2/2\Dt}$. Since
$\nabla_x g_{\Dt}(x-e^{-t}a) = -\frac{x-e^{-t}a}{\Dt}g_{\Dt}(x-e^{-t}a)$ and, for $x$ in a
compact neighbourhood, $\|x-e^{-t}a\|\,g_{\Dt}(x-e^{-t}a) \le c_1 e^{-c_2\|a\|^2}$ uniformly, the
integrand is dominated by an integrable function whenever $\Pz$ has a finite first moment.
Dominated convergence then licenses
\[
\nabla_x \Pt(x)
= \int \mathrm{d}a\; \Pz(a)\,\frac{e^{-t}a - x}{\Dt}\, g_{\Dt}(x-e^{-t}a).
\]
Dividing by $\Pt(x)$ and recognising $\Pz(a)g_{\Dt}(x-e^{-t}a)/\Pt(x) = P(a\mid x,t)$ gives
\eqref{eq:score}. Nothing is approximated. In the empirical-Bayes literature this is Tweedie's
formula \citep{robbins1956,miyasawa1961,efron2011}.

\paragraph{Risk decomposition.} For any $u$ that is a function of $x$ alone,
\begin{equation}
\label{eq:orthogonal}
\E\!\left[\|a-u\|^2 \,\middle|\, x\right]
= \E\!\left[\|a - \pmean{a}\|^2 \,\middle|\, x\right] + \|u - \pmean{a}\|^2 ,
\end{equation}
by expanding $a-u = (a-\pmean{a}) + (\pmean{a}-u)$ and noting the cross term vanishes because
$\pmean{a}-u$ is $x$-measurable and $\E[a-\pmean{a}\mid x]=0$. Hence $\pmean{a}$ minimises the
conditional risk, with an irreducible floor given by the first term \citep{lehmann1998}.
Optimality is with respect to squared error specifically; a different loss selects a different
functional of the posterior. This is why every table in \S\ref{sec:networks} reports error
\emph{against} $\pmean{a}$ rather than raw risk: the floor is common to all estimators and
subtracting it is what makes the comparison informative.

\section{What ``exact'' means here, in four parts}
\label{app:exactness}

The body compresses this to two sentences; the distinctions matter enough to be kept apart
somewhere, because collapsing any two of them is how a discretised number acquires the authority
of an exact one.

(i) Sum-product on \eqref{eq:posterior} is \emph{exact at the level of functional messages}, as a
consequence of the graph being a tree.

(ii) A message is a \emph{function on $\Reals$}, and this is forced by the state space, not chosen.
On a discrete graphical model $a_i$ takes finitely many values and $\lmsg_i$ is a vector of that
length. Here $a_i$ ranges over $\Reals$, so $\lmsg_i$ assigns a number to every real $a_i$: it is
an unnormalised density, the evidence from everything left of site $i$. Only when the transition
and the likelihood are both Gaussian does that function stay inside a finite-dimensional family
--- two numbers, a mean and a variance, which is the Kalman filter. For a general $\innov$ it does
not, and the integral in \eqref{eq:messages} produces a new shape at every site. An implementation
must therefore represent a function. We use a truncated uniform grid
$[-\halfwidth,\halfwidth]$ with $\ngrid$ points, spacing $\gridh$ and composite trapezoidal
weights, which defines a \emph{finite-state model} for which the recursion is \emph{exact up to
floating-point arithmetic}.

(iii) That finite-state model \emph{approximates} the continuous one, with an error that must be
measured --- \S\ref{sec:rung0} measures it.

(iv) Monte Carlo enters this note in exactly one place: when a statistic is computed from
\emph{generated samples}. Every other number here --- posterior means, evidences, information
matrices --- comes from deterministic quadrature and carries no sampling error at all, only
discretisation error. Where a quantity is averaged over random draws of a dataset, that is
replicate spread and is labelled as such.

Three levels of exactness are involved and collapsing any two of them overstates the result, so
each has its own name and none is abbreviated to ``exact''.

\begin{enumerate}
\item \emph{Tree-exact functional BP.} The posterior factorises on a chain and the sum-product
integrals are exact as statements about functions on $\Reals$. Nothing numerical has happened yet.
\item \emph{Exact recursion for the discretised model.} Once grid and quadrature weights are
fixed, forward--backward is exact --- up to floating point --- \emph{for that finite-state model}.
It is not exact for the continuous chain.
\item \emph{Validated grid-quadrature reference.} The finite-state model approximates the
continuous chain, with truncation and quadrature error that is measured rather than assumed
(\S\ref{sec:rung0}).
\end{enumerate}

Accordingly the tables say ``deterministic grid-BP reference under the true kernel'', and the
population statement is that the score \emph{reduces to tree-exact functional inference and is
evaluated numerically by controlled grid quadrature}. ``Exact grid-BP'' would collapse (ii) into
(iii) and quietly promote a measured approximation into a guarantee.

\section{Sum-product on the posterior chain}
\label{app:bp}

\paragraph{Posterior factorisation.} Bayes' rule applied to \eqref{eq:prior}, with the likelihood
factorising sitewise because the forward noise is independent across coordinates, gives
\eqref{eq:posterior}. Each likelihood factor has degree one among latent variables, so attaching
it to the prior's path adds a leaf and cannot create a cycle: the posterior factor graph is a
tree. Conditioning usually creates dependence; it does not here, precisely because each
observation is generated from a single latent variable. The construction fails the moment the
forward noise is correlated across coordinates, since then $P(x\mid a)$ does not factorise, the
likelihood becomes a factor touching several $a_i$, and the posterior graph acquires edges.

\begin{figure}[!htbp]
\centering
\begin{tikzpicture}[
  latent/.style={circle, draw, minimum size=7mm, inner sep=1pt},
  obs/.style={circle, draw, fill=black!10, minimum size=7mm, inner sep=1pt},
  fac/.style={rectangle, draw, fill=black, minimum size=2.2mm, inner sep=0pt},
  node distance=9mm]
\foreach \i in {1,2,3,4}{
  \node[latent] (a\i) at (2.1*\i, 0) {$a_{\i}$};
  \node[fac] (f\i) at (2.1*\i, -1.0) {};
  \node[obs] (x\i) at (2.1*\i, -2.0) {$x_{\i}$};
  \draw (a\i) -- (f\i) -- (x\i);
}
\foreach \i/\j in {1/2,2/3,3/4}{
  \node[fac] (w\i) at (2.1*\i+1.05, 0) {};
  \draw (a\i) -- (w\i) -- (a\j);
}
% The transition-factor label sits just above its factor square; the two
% message arrows are lifted clear of it so nothing overlaps.
\node[above=1mm of w1] {\small $\kernel(a_{i}\mid a_{i-1})$};
\node[right=2mm of f4] {\small $\like(x_i\mid a_i)$};
\draw[-{Stealth}, thick] (1.15,1.35) -- (2.55,1.35) node[midway,above] {\small $\lmsg$};
\draw[-{Stealth}, thick] (9.35,1.35) -- (7.95,1.35) node[midway,above] {\small $\rmsg$};
\end{tikzpicture}
\caption{Factor graph of the posterior \eqref{eq:posterior}. Open circles are latent variables,
shaded circles are observations, squares are factors. Each transition factor
$\kernel(a_i\mid a_{i-1})$ joins two adjacent latent variables; each observation enters through a
\emph{unary} factor $\like$ attached to a single latent variable. Conditioning therefore reweights
node potentials without creating edges among the $a_i$, and the posterior graph remains a chain.
Arrows show the two message sweeps \eqref{eq:messages}.}
\label{fig:factorgraph}
\end{figure}

\paragraph{Exactness on a tree.} Let the graph be a tree and let $a_i$ be a variable node.
Removing $a_i$ disconnects it into components, and every factor involves $a_i$ together with
variables of exactly one component --- a factor touching two components would create a second
path between them, hence a cycle. Grouping factors by component gives
$P \propto \prod_c g_c(a_i, v^{(c)})$, so conditionally on $a_i$ the components are independent.
Unrolling the forward recursion \eqref{eq:messages},
\[
\lmsg_i(a_i) \propto \int\!\mathrm{d}a_{1:i-1}\;
\mu(a_1)\prod_{j=2}^{i}\kernel(a_j\mid a_{j-1})\prod_{j=1}^{i-1}\like(x_j\mid a_j),
\]
which is the joint density of everything at or left of site $i$ with the left variables
integrated out and $\like(x_i\mid a_i)$ \emph{not} included; symmetrically for $\rmsg_i$. The two
are conditionally independent given $a_i$, so their product times the local evidence is
proportional to the marginal. The pairwise statement is the same argument with the separator
taken to be the edge.

\paragraph{Where the local factor goes.} In \eqref{eq:messages} the local evidence at site $i-1$
is absorbed \emph{before} the transition, so $\lmsg_i$ carries the factors
$\like(x_1\mid a_1),\dots,\like(x_{i-1}\mid a_{i-1})$ and $\rmsg_i$ carries
$\like(x_{i+1}\mid a_{i+1}),\dots,\like(x_{\nsites}\mid a_{\nsites})$, and each appears exactly
once in the marginal. Absorbing it \emph{after} would place $\like(x_i\mid a_i)$ in both
messages and square it: the
belief would be sharper than the truth and every posterior variance too small. This produces
plausible-looking output, which is why the convention is pinned by
\codefile{tests/test\_message\_normalization.py} rather than left to a comment.

\paragraph{Evidence.} Messages are renormalised at each step to prevent underflow; accumulating
the discarded constants recovers $\log\Pt(x)$ for the discretised model exactly. This is the
objective that \S\ref{sec:rung1} maximises, and having it in closed form is what makes it possible
to check a gradient against a finite difference of the thing being optimised.
\coderef{src/bp\_grid.py}{grid\_bp} returns it alongside the beliefs;
\coderef{src/exact\_scores.py}{gaussian\_log\_evidence} is the closed form it is tested against.

\paragraph{Complexity at the functional level.} Each message update integrates one variable
against the kernel. Abstractly this is an integral operator applied to a function; represented on
$\ngrid$ points it is one $\ngrid\times\ngrid$ matrix--vector product, $O(\ngrid^2)$. There are
$2\nsites$ updates, so one sequence costs $O(\nsites\ngrid^2)$ against $O(\ngrid^{\nsites})$ for
naive marginalisation. At $\nsites=32,\ngrid=401$ that is about $5\times10^6$ operations.
Batching over sequences turns each update into a single GEMM, which is why the same code runs on
a GPU without restructuring (\codefile{src/backend.py}, gated by
\codefile{tests/test\_backend\_parity.py}).

\section{Covariance-stationary versus strictly stationary}
\label{app:stationary}

Two notions of stationarity coincide for Gaussian processes and separate here, and the body
relies on the weaker one, so the difference is worth stating.

A process is \emph{strictly stationary} if the joint law of $(a_{i+1},\dots,a_{i+k})$ does not
depend on $i$, for every $k$. It is \emph{covariance-stationary} (or weakly, or second-order
stationary) if only the first two moments are shift-invariant: $\E[a_i]$ constant and
$\Cov(a_i,a_j)$ a function of $|i-j|$ alone. Strict stationarity implies the covariance version
whenever second moments exist; the converse is false in general.

\paragraph{Why they coincide in the Gaussian case.} A Gaussian law is determined by its mean
vector and covariance matrix. If the mean is constant and the covariance depends only on $|i-j|$,
then the joint law of any window $(a_{i+1},\dots,a_{i+k})$ is the Gaussian with that mean and that
covariance block --- and the block depends on $i$ only through the differences $|j-j'|$, which are
shift-invariant. So the joint law is shift-invariant, which is strict stationarity. The
implication runs through ``determined by two moments'' and therefore fails for any family that is
not.

\paragraph{What holds for our chain.} With $\ivar=1-\corr^2$ and $\Var(a_1)=1$, second moments
propagate exactly:
\[
\Var(a_i) = \corr^2\Var(a_{i-1}) + \ivar = \corr^2 + (1-\corr^2) = 1 ,
\]
by induction, and $\Cov(a_i,a_j)=\corr^{|i-j|}$, \emph{whatever the shape of $\innov$}. The chain
is therefore covariance-stationary, and this is exactly the property the variance-matched family
needs: Gaussian, Laplace, Student, uniform and bimodal innovations share the same second-order
structure and differ only beyond it.

It is \textbf{not} strictly stationary for non-Gaussian $\innov$. The invariant law of
$a_i=\corr a_{i-1}+\varepsilon_i$ is that of $\sum_{k\ge0}\corr^k\varepsilon_{-k}$, an infinite
convolution which is not Gaussian, so drawing $a_1\sim\normal{0}{1}$ starts the chain off its
invariant law. The marginal of $a_i$ then drifts towards it over a burn-in of order
$1/\log(1/\corr)\approx 6$ sites at $\corr=0.85$. This does not affect the comparisons in the
body, because every family shares the same $a_1$ law and the same covariance and so still differs
only beyond second moments --- but $\ivar=1-\corr^2$ alone does not deliver strict stationarity,
and we do not claim it does.

\section{The Gaussian model and the second-order baseline}
\label{app:lmmse}

For $\innov=\normal{0}{\ivar}$ with $\ivar=1-\corr^2$ and $a_1\sim\normal{0}{1}$, the vector $a$
is jointly Gaussian with $(\Sigma_0)_{ij}=\corr^{|i-j|}$. Expanding
$\Pz(a)\propto\exp[-\tfrac{1}{2\ivar}\sum_{i\ge2}(a_i-\corr a_{i-1})^2 - \tfrac12 a_1^2]$
produces only $a_i^2$ and $a_ia_{i-1}$ terms, so $\Sigma_0^{-1}$ is tridiagonal with
$(\Sigma_0^{-1})_{ii} = (1+\corr^2)/\ivar$ in the interior, $1/\ivar$ at the two ends, and
$-\corr/\ivar$ on the off-diagonals. Adding the likelihood contributes $(e^{-2t}/\Dt)I$, so the
posterior precision is tridiagonal at every $t$: Markovianity is a statement about the precision,
while the covariance $\Sigma_0$ is dense.

Since $x = e^{-t}a + \sqrt{\Dt}z$ with $z$ independent, $(a,x)$ is jointly Gaussian with
$\Cov(a,x) = e^{-t}\Sigma_0$ and $\Cov(x) = e^{-2t}\Sigma_0 + \Dt I =: \Sigma_t$, whence
$\E[a\mid x] = e^{-t}\Sigma_0\Sigma_t^{-1}x$ and, by \eqref{eq:score},
$\score(x,t) = -\Sigma_t^{-1}x$. For zero-mean jointly Gaussian variables the conditional
expectation is linear, so it coincides with the LMMSE estimator \citep{kay1993}. That supplies the
covariance identity the proposition's proof appeals to; the proof itself, an induction over the
forward and backward recursions establishing that every projected message equals the Gaussian
chain's, follows.

%% SHARED SECTION -- the proof of Proposition~\ref{prop:lmmse}, included from
%% the appendix. Split out of prop-lmmse.tex because the body has a nine-page
%% limit and a full induction does not belong in it; the STATEMENT, the
%% projection definition and the hypothesis discussion stay in the body, which
%% is what a reader needs there.

\begin{remark}[On the hypotheses]
\label{rem:lmmse-hyp}
Four are load-bearing and worth separating. \emph{Finite second moments} are what make $\Pgauss$
defined at all. \emph{Zero mean} is not cosmetic: the general LMMSE carries
$\E[\va] + \Cov(\va,\vx)\Cov(\vx)^{-1}(\vx-\E[\vx])$, and \eqref{eq:lmmse} is the centred
special case. \emph{$\corr\neq0$} is required because the backward map in \eqref{eq:closure}
divides by $\corr$; at $\corr=0$ the chain has no edges, the sites are independent, and the
claim holds trivially and separately, each site's posterior mean being the scalar Gaussian
$\Sig_0(\Sig_0+\Dt e^{2t})^{-1}x_i$. \emph{Where the projection is applied} matters because
$\Pgauss$ does not commute with multiplication by the likelihood; Proposition~\ref{prop:lmmse}
fixes the order, and a different order is a different algorithm.

The statement is about exact message functions on $\Reals$ and an analytic projection operator.
\textbf{The grid plays no part in it.} Discretisation enters only when the \emph{unrestricted}
non-Gaussian messages are represented numerically, which is what the recursion is compared
against; conflating the two would turn a theorem about a closure into a claim about quadrature.
\end{remark}

\begin{proof}
Write $\mu_i^{\rightarrow}$ for the forward message at site $i$ after projection and
$\ell_i$ for the Gaussian unary likelihood. We show by induction that every projected message
coincides with the corresponding exact message of the covariance-matched Gaussian chain --- the
chain with the same $\corr$ and innovation variance $\ivar$ but Gaussian innovations.

\emph{Base.} The initial message is Gaussian by hypothesis, or is projected to the Gaussian with
the same first two moments; either way it equals the Gaussian chain's initial message, which has
those same moments by construction.

\emph{Step.} Suppose $\mu_{i-1}^{\rightarrow}$ agrees with the Gaussian chain's message, so it is
Gaussian with some mean $m$ and variance $v$. Multiplying by $\ell_{i-1}$ multiplies two Gaussian
densities and yields a Gaussian, with mean and variance given by the usual precision-weighted
combination --- identical in both chains, since the likelihood is the same object. Taking the
transition integral against $a_i = \corr a_{i-1} + \varepsilon_i$ convolves with $\innov$ and
scales by $\corr$, which maps $(m,v) \mapsto (\corr m,\, \corr^2 v + \ivar)$: the mean and variance
of a convolution depend on $\innov$ \emph{only} through its mean and variance, so this is the same
map for every innovation with moments $(0,\ivar)$, and in particular the same as for the Gaussian
chain. In the Gaussian chain the convolution of two Gaussians is already Gaussian, so no
projection occurs there; here $\Pgauss$ is applied, and since the pre-projection function has
exactly the moments just computed, the projection returns precisely the Gaussian chain's message.
The backward recursion is identical with the transition map inverted, $(m,v)\mapsto
(m/\corr,\,(v+\ivar)/\corr^2)$, which is where $\corr\neq0$ is used.

Hence every projected belief equals the corresponding Gaussian-chain belief, and the returned
posterior mean is the Gaussian chain's posterior mean. For jointly Gaussian zero-mean
$(\va,\vx)$ the conditional mean is linear in $\vx$ and therefore coincides with the best linear
estimator, which is \eqref{eq:lmmse}: since $\vx = e^{-t}\va + \sqrt{\Dt}\,\vect{z}$ with
$\vect{z}$ independent, $\Cov(\va,\vx)=e^{-t}\Sig_0$ and $\Cov(\vx)=e^{-2t}\Sig_0+\Dt\Id$. Since
the recursion never read $\innov$ beyond $(0,\ivar)$, the same estimator is returned for every
innovation law with those moments.
\end{proof}


Consequently, on this model, the \emph{message-representation error} and the \emph{model error}
coincide: projecting messages onto Gaussians and replacing the prior by its covariance-matched
Gaussian give the same estimator. This is a property of single-Gaussian closure with linear
transitions, not of Gaussian approximations in general, and the two separate for richer message
families.

\section{Numerical representation}
\label{app:numerics}

Messages are represented on a uniform grid $\gridpt{1},\dots,\gridpt{\ngrid}$ spanning
$[-\halfwidth,\halfwidth]$ with spacing $\gridh = 2\halfwidth/(\ngrid-1)$ and composite
trapezoidal weights $w_k = \gridh$ for interior points, $\gridh/2$ at the two ends.

The recursions are \emph{scaled linear-domain} forward--backward, not log-semiring belief
propagation. The transition matrix and the unary likelihoods are built in logs --- the likelihood
with a per-site maximum subtracted before exponentiating --- and are then exponentiated once, so
each message update is an ordinary $(\ngrid,\ngrid)$ matrix product followed by renormalisation to
unit quadrature mass. The removed constants are accumulated and restored in the evidence, which is
therefore exact. This is the standard scaled recursion and it is what keeps the \emph{total} mass
of every message at unity, so no message vanishes as a whole.

It does not, however, guarantee that no individual entry underflows: a single likelihood or
transition value far in a tail can round to zero while the message it belongs to remains
well-normalised. The recursion is scaled, not logarithmic, so this is a real possibility and is
measured below rather than excluded by construction. The $10^{-14}$ agreement reported below is
correspondingly evidence at the Gaussian configuration tested, not a uniform bound over innovation
families, noise levels and fitted kernels.

\paragraph{How far from a bound, measured.} Sweeping $\ngrid\in\{201,401,801\}$ and
$\halfwidth\in\{6,8,10\}$ against a finer reference ($\ngrid=1601$, $\halfwidth=12$), over all
twelve noise levels and four innovation families, with observations drawn from the prior and again
conditioned into the tail at ${\sim}3$ standard deviations:

\begin{center}\small
\begin{tabular}{@{}lcccc@{}}
\toprule
worst relative score error at $\ngrid=401,\ \halfwidth=8$
  & Gaussian & bimodal & Laplace & Student \\
\midrule
observations drawn from the prior & $2\times10^{-13}$ & $1\times10^{-15}$
  & $1\times10^{-5}$ & $\mathbf{9.9\times10^{-4}}$ \\
observations conditioned into the tail & $2\times10^{-7}$ & $1\times10^{-11}$
  & $\mathbf{1.7\times10^{-2}}$ & $8.3\times10^{-3}$ \\
\bottomrule
\end{tabular}
\end{center}

\noindent The $10^{-14}$ figure is the \emph{best} cell in this table, not a bound on it. Under
ordinary draws the production grid resolves the whole schedule to $10^{-3}$ or better across every
family, which is the regime all reported results are in and is the threshold the generator
enforces. Under displaced observations the two heavy-tailed families lose two further orders.
Partial underflow is real rather than hypothetical: up to $0.27\%$ of likelihood entries round to
exactly zero at the production configuration, rising to $3.8\%$ at $\ngrid=201$ --- always under
displacement, never under ordinary draws.

\paragraph{What the stress row is, exactly.} It is a three-standard-deviation
\emph{displacement} of the whole chain, not conditioning on a tail event of the prior. A translate
and a conditional law are different objects, and it is the translate that is computed: a
legitimate stress test, named as one.

\paragraph{Boundary occupancy, measured inside the recursion.} The quantity that matters is how
much mass the objects BP actually propagates place at the edge of the domain, so it is measured on
those objects rather than on a proxy assembled from the likelihood and the transition matrix. The
worst forward-message edge occupancy over the schedule is $\gcfwdedge$ and the worst node-belief
edge occupancy is $\gcbeledge$, both over $\gccells$ production cells. The worst per-site
log-evidence shift against the reference resolution is $\gcworstev$ nats.

\paragraph{Fitted kernels, the hardest case for the grid.} The families above are four analytic
priors, but the kernels the reported runs actually use are \emph{fitted} mixtures, which can place
a component far narrower than any smooth law. Since the narrowest component relative to the mesh is
precisely what the grid must resolve, this is the case most likely to fail and is tested directly.
Adding a fitted mixture with a component two grid cells wide changes the picture in one place and
not the other. Under ordinary draws it is unremarkable and the schedule still
resolves to $\gcworstord$. Under three-standard-deviation displacement at the narrowest domain
tried ($\halfwidth=6$) it is the worst cell in the entire sweep, with $96\%$ of node-belief mass
against the boundary: a narrow component displaced toward the edge of a too-small window is not
resolved at all. That configuration is outside the production setting, and it marks where the
method stops working rather than a defect in the reported numbers --- but it is the failure mode to
check first if the domain is ever narrowed.

\paragraph{What this does and does not certify.} It is a self-convergence stress test on fresh
synthetic draws. It establishes that errors on representative ordinary draws are below $10^{-3}$.
It does \emph{not} establish that no reported number moves: that would require re-evaluating the
actual held-out observations and fitted kernels at a finer resolution --- a different experiment,
and one not run here. Section~\ref{app:limits} records the residual limitation.

\begin{figure}[h]
\centering
\includegraphics[width=0.72\textwidth]{fig_grid_domain.pdf}
\caption{Grid and domain control, Laplace chain. (a) Relative score error against a
high-resolution reference at fixed $\halfwidth=8$; (b) the same at fixed $\ngrid=401$ as the
domain grows. Both saturate at the floating-point floor $\sim10^{-14}$, so beyond that the curves
measure arithmetic, not discretisation. (c) Forward-message boundary mass against $\halfwidth$ ---
the truncation diagnostic --- as the worst case and the $90$th percentile over $256$ chains, which
differ by roughly two orders of magnitude and are therefore both drawn. (d) Interior column-mass
residual against spacing at $\halfwidth=8$ --- the quadrature diagnostic --- with an $O(\gridh^2)$
reference. The two diagnostics do different jobs and neither alone certifies convergence.}
\label{fig:grid}
\end{figure}

\paragraph{Two error sources.} \emph{Truncation}: probability mass outside
$[-\halfwidth,\halfwidth]$ is discarded rather than transported, an error controlled by
$\halfwidth$ and essentially independent of $\ngrid$. \emph{Quadrature}: the trapezoidal rule is
second order in $\gridh$ for integrands with bounded second derivative, an error controlled by
$\gridh$ at fixed $\halfwidth$. Figure~\ref{fig:grid} measures them separately, using the
forward-message boundary mass for the first and the interior column-mass residual of the
transition matrix for the second; the interior restriction $|\gridpt{}| \le \halfwidth/2$ removes
columns whose transition density has a tail outside the grid, which is what makes the second
diagnostic measure quadrature alone.

The boundary mass is the maximum over sites of a \emph{sampled} trajectory, so it is a statistic
and not a bound, and it is reported over $256$ chains per cell rather than from one draw. At
$\ngrid=401$, taking the worst case over the twelve-level schedule,

\begin{center}\small
\begin{tabular}{@{}lcccc@{}}
\toprule
$\halfwidth$ & $4$ & $6$ & $8$ & $10$ \\
\midrule
worst chain & $2.3\times10^{-3}$ & $1.6\times10^{-4}$ & $1.2\times10^{-6}$ & $3.3\times10^{-8}$ \\
$90$th percentile & $2.7\times10^{-4}$ & $2.5\times10^{-6}$ & $1.5\times10^{-8}$ & $7.8\times10^{-11}$ \\
median & $3.8\times10^{-5}$ & $3.8\times10^{-7}$ & $3.0\times10^{-9}$ & $1.8\times10^{-11}$ \\
\bottomrule
\end{tabular}
\end{center}

\noindent
The spread between the rows is the point: at the working $\halfwidth=8$ the worst chain carries
$84$ times the boundary mass of the $90$th percentile, so a single trajectory's value --- or an
upper quantile quoted as though it were a bound --- understates the tail by close to two orders of
magnitude. The interior residual falls by a factor of four per halving of $\gridh$, i.e.\
$O(\gridh^2)$.

At $\halfwidth=4$ the interior residual saturates at $\sim10^{-3}$ instead of continuing to fall,
because at that domain even the ``interior'' columns have truncated tails --- the floor is the
truncation error leaking into the quadrature diagnostic. This is the reason the two diagnostics
have to be read together and neither alone certifies convergence.

\paragraph{What is not claimed.} A grid-resolution sweep at fixed $\halfwidth$ does not rule out
truncation error, and a domain sweep at fixed $\ngrid$ does not rule out quadrature error. We
report both axes and the two diagnostics; we do not claim continuum convergence, only that at
$\ngrid=401$, $\halfwidth=8$ both diagnostics sit far below every effect the paper reports.
For the Laplace family, whose innovation density has a discontinuous derivative at the origin, we
claim second order and no better.

\section{Learning procedure}
\label{app:learning}

\paragraph{Objective.} Maximum marginal likelihood
$\Lik(\params) = \sum_\mu\log P_{t,\params}(x^\mu)$, with the complete-data objective
$\log P_\params(a,x) = \log\mu(a_1) + \sum_{i\ge2}\log\kernelth(a_i\mid a_{i-1}) +
\sum_i\log\like(x_i\mid a_i)$. Only the middle term depends on $\params$, which is why $\Xis$ of
\eqref{eq:xi} is sufficient and why $\Qfun(\params\mid\params') = \langle\Xis(\params'),
\log\kernelth\rangle$ up to $\params$-free constants.

\paragraph{E-step.} One forward and one backward sweep per sequence, then a rank-one accumulation
of the pairwise statistics: writing $\ell^{(i)}$ for the vector $\like(x_i\mid \cdot)$ on the
grid, and with $f_i = w\odot\lmsg_i\odot\ell^{(i)}$ and $g_i = w\odot\rmsg_i\odot\ell^{(i)}$,
the pairwise belief on edge $(i,i+1)$ is
$f_i(j)\,\kernel[k,j]\,g_{i+1}(k)/Z_i$ with $Z_i = g_{i+1}^\top \kernel f_i$, and
$\Xis = \left(\sum_{\mu,i} g_{i+1}(f_i)^\top/Z_i\right)\odot \kernel$. $\Xis$ sums to the number
of edges in the dataset, which is asserted in the tests. Exact for the discretised model, in the
sense of Appendix~\ref{app:exactness}(ii). Implementation \coderef{src/em.py}{e\_step}.

\paragraph{M-step, by kernel.}
\begin{itemize}\itemsep2pt
\item \emph{Gaussian}, $\params=(\corr,\ivar)$: closed-form maximiser, the belief-weighted
      Yule--Walker equations $\corr = S_{01}/S_{00}$, $\ivar = (S_{11}-\corr S_{01})/W$. Exact
      maximisation, hence exact EM for that family.
\item \emph{Laplace}, $\params=(\corr,b)$: closed form despite non-differentiability. Profiling
      $b$ out of $\Qfun$ turns the problem into minimising the belief-weighted mean absolute
      residual, whose exact solution is a weighted median. Exact maximisation.
\item \emph{Mixture} \eqref{eq:kernel}: the complete-data problem carries a second latent
      variable, so the M-step is itself an inner EM over the mixture label with
      $\Xis$ held fixed. We alternate the $(\pi,\nu,s^2)$ block and the $\corr$ block, both
      closed form, for at most \textbf{sixteen} inner sweeps, stopping early when the per-edge
      gain in $\Qfun$ drops below $10^{-8}$. This increases $\Qfun$
      without maximising it: \textbf{generalised EM / ECM} \citep{meng1993}, not exact EM.
      Sixteen is chosen by measurement: at four sweeps the fitted innovation excess kurtosis
      is $28\%$ below its value at sixteen whenever the outer loop is also cut short, while
      $\corr$ moves in the fourth decimal --- the same shape-versus-correlation rate asymmetry
      the outer loop shows, one level down. Since shape is the coordinate every non-Gaussian
      claim depends on, the budget is reported here rather than left as a default.
      Because each block is a conditional \emph{maximiser}, $\Qfun$ cannot decrease along an
      inner sweep; the implementation raises rather than stopping if it ever does, and across
      the certified fits it never has.
\item \emph{Mixture-density network}: one backtracked Adam step, accepted only if $\Qfun$
      increased, halving the step up to five times otherwise and standing still if no ascent is
      found. Also generalised EM.
\end{itemize}

\paragraph{Guarantees, initialisation, stopping.} Generalised EM guarantees monotone ascent of
$\Lik$ \citep{dempster1977,wu1983}; convergence is to a stationary point, not necessarily a global
maximum. Mixture components are initialised spread over $[-\sqrt{\Var}, \sqrt{\Var}]$ with random
scales; $\corr$ is initialised away from its truth in every experiment --- at $0.3$ in most, and
over $\{0,\,0.3,\,0.6,\,-0.4\}$ in the initialisation sweep of \codefile{exp\_18}. The truth is
$0.85$ in \codefile{exp\_02}, \codefile{exp\_16}, \codefile{exp\_18} and \codefile{exp\_22}, and
$0.8$ in \codefile{exp\_06}, \codefile{exp\_07} and \codefile{exp\_08}; every table and figure
states the value it used. Iteration
stops when the relative increase of $\Lik$ falls below $10^{-9}$, capped at $50$--$120$
iterations depending on the experiment. Variances are floored at $10^{-6}$ to prevent component
collapse; no other regularisation is applied, and the component means are unconstrained
(\S\ref{sec:rung1}). Monotonicity is recorded per run as the largest decrease across iterations and
is zero in every run reported here.

\section{Convergence rates, clean against channel}
\label{app:rates}

Two slownesses, measured separately, because they are independent and easily conflated.

\paragraph{Optimisation: coordinates of the same kernel settle at very different rates.} Running
EM to $800$ updates at $\nseq=512$, $\ncomp=8$ and recording the whole parameter trace, we measure
for each coordinate the last update at which it was still further than a fixed relative tolerance
from its final value. Over sixteen seeds the median is $80$ updates for $\corr$ ($10^{-3}$
tolerance), $68$ for the innovation variance ($10^{-2}$), and $229$ for the excess kurtosis
($2\times10^{-2}$), the last ranging up to $638$. The shape is thus roughly three times slower to
settle than the correlation, and monitoring $\corr$ --- the natural thing to plot --- shows a
converged-looking trace over a kernel that is not. This is why the stopping rule may not be read
off $\corr$, and why an iteration budget must be justified against the \emph{slowest} coordinate.

\paragraph{Statistics: fitting through the channel costs rate, not just constant.} Comparing
maximum likelihood on clean transition pairs against the same estimate obtained through the
noising channel, at sixteen replicates over $\nseq=64$ to $1024$, the fitted $\log$--$\log$ slopes
are

\begin{center}\small
\begin{tabular}{@{}lcc@{}}
\toprule
& $\corr$ & innovation variance \\
\midrule
clean transition pairs & $-0.52$ & $-0.50$ \\
through the channel    & $-0.44$ & $-0.36$ \\
\bottomrule
\end{tabular}
\end{center}

\noindent
The clean arm sits on the parametric $\nseq^{-1/2}$ to two decimal places on both coordinates,
which is the check that the estimator itself is not the bottleneck; the noised arm is slower, and
more so on the shape than on the correlation. Note that the clean arm is \emph{not} flat in
$\nseq$ --- its error falls by a factor of $3.4$ on $\corr$ and $4.3$ on the variance across the
range. This is worth stating explicitly: a flat clean arm would indicate a discretisation floor,
and at three replicates the curve is noisy enough to look flat when it is not.

\begin{figure}[t]
\centering
\includegraphics[width=0.55\textwidth]{fig_innovation_density.pdf}
\caption{The fitted innovation density against the true Laplace density, on a log scale so the
tails are visible. What is recovered here is a density, not a parameter: the mixture is a
representation of it.}
\label{fig:density}
\end{figure}

\section{The fitted rates in the efficiency comparison}
\label{app:fitted-rates}

The body reports fitted $\log$--$\log$ slopes of relative denoising error against
$\nseq$. They are fitted \emph{per seed} and then averaged, not fitted to the
seed-averaged curve: cells within a seed share one training set and one set of
fitted models, so the two differ, and only the first keeps the pairing that makes
the arms' difference testable. Over $\nseq=32$ to $\biggestn$ the network gives
$\ratenet \pm \ratenetse$ and EM--BP $\rateem \pm \rateemse$, a paired difference
of $\rategap \pm \rategapse$ ($t = \rategapt$ on $\rateseeds - 1$ degrees of
freedom). The sign test is the one to trust at this sample size, and it is
unanimous: EM--BP is steeper in all $\ratesteeperseeds$ seeds. Restricting to
$\nseq \ge 256$ or to $\nseq \le 512$ leaves the difference at $0.13$ and $0.10$
respectively, so it is not carried by either end of the range.

\paragraph{Why censoring does not explain it.} Both arms select their largest
allowed budget in a growing fraction of cells as $\nseq$ grows --- the network
from $6\%$ at $\nseq=32$ to $42\%$ at $\biggestn$, EM--BP from $12\%$ to $61\%$.
By the rule used throughout this work, two arms with different cap-hit rates are
being compared on their budgets, so the fitted slopes are uninterpretable unless
the direction of the effect is known. It is measured. Tripling both caps at
$\nseq=\nseqcalib$, on the same seeds and validation bundles, lowers the network's
error by $\ratecalibnet\%$ and EM--BP's by $\ratecalibem\%$ --- an order of
magnitude apart. EM--BP is the arm being held back, and it is also the arm whose
cap-hit rate rises faster with $\nseq$, so the censoring inflates its error more
at large $\nseq$ than at small and \emph{flattens} its fitted curve. Removing the
censoring would steepen EM--BP's slope and widen the difference. The reported
value is a lower bound.

Two further checks point the same way. Fitting the slope between $\nseq=\nseqcalib$
and $\nseq=\biggestn$ using only the raised-budget runs at both ends --- so that
neither endpoint is at the certified cap --- gives a difference of $0.26 \pm 0.12$,
larger than the full-range figure, though a two-point slope on $\rateseeds$ seeds
is individually weak. And the fit is generated rather than transcribed: the
generator refuses to emit if the budget calibration ever shows the network as the
more censored arm, since that is the single assumption the claim rests on which
the fitted slopes do not themselves display.

\paragraph{What the slopes are not.} Neither matches a rate we can derive.
EM--BP's $\rateem$ is well short of the $\nseq^{-1/2}$ that
Appendix~\ref{app:protocol} and the Rung~1 measurement obtain for $\ivar$ alone,
which is consistent with a pointwise error inheriting the slower of its two
coordinates rather than the faster --- $\corr$ falls at $-0.27$ over the same
range --- but we have not established that, and we do not build on it. The claim
in the body is the qualitative one.

\section{Experimental protocol}
\label{app:protocol}

%% SHARED SECTION -- the scope remark for Proposition~\ref{prop:fisher},
%% included from the appendix. The body carries the statement; a reader asking
%% "exact for WHICH likelihood?" is served here, where there is room to answer.

\begin{remark}[Discrete and continuous are separate claims]
\label{rem:fisher-scope}
Proposition~\ref{prop:fisher} is a statement about the \emph{finite-state model on the grid}. There
the interchange of $\nabla_\params$ with a finite sum is unconditional, and the returned gradient
is exact for the discretised likelihood up to floating point --- which is what the
$10^{-9}$ finite-difference agreement of Appendix~\ref{app:protocol} measures.

The corresponding continuum identity,
$\nabla_\params \Lik = \int\!\!\int \xi_\params(a,a')\,\nabla_\params \log
\kernel_\params(a'\mid a)\,\dd a\,\dd a'$, also holds, but it needs hypotheses the discrete
statement does not: $\kernel_\params(a'\mid a)$ differentiable in $\params$, support not depending
on $\params$, and an integrable dominating function for $\partial_\params \kernel_\params$ so that
differentiation and integration may be exchanged.

Gaussian and Gaussian-mixture innovations satisfy the hypotheses \emph{as written}: both are
smooth in $\params$ with full support. Laplace does not, and saying so matters because Laplace is
the innovation law the reported runs use. Its density $\propto \exp(-|a'-\corr a|/b)$ has a kink
at $a'=\corr a$, so $\partial_\corr \kernel_\corr$ fails to exist on that set and the proposition's
smoothness hypothesis is violated there. The identity survives, but by a different argument: the
kink set is Lebesgue-null, the derivative exists almost everywhere, and the difference quotients
are dominated locally uniformly in $\corr$, so dominated convergence still licenses the exchange.
The discrete statement is untouched either way --- on the grid the kink is a finite set of entries
and the sum is differentiated term by term --- and it is the discrete one the reported numbers
certify.
\end{remark}


\paragraph{What $\Xis$ is sufficient for.} The complete-data log-likelihood is a sum over edges, so
the $\params$-dependent part of $\Qfun(\params\mid\params^{(r)})$ sees the data only through
$\Xis$. But $\Xis$ is itself computed at $\params^{(r)}$ and rebuilt every iteration, so the claim
concerns one maximisation step, not the data being compressible once. It also requires a
homogeneous kernel shared across sites and sequences, a fixed grid, an initial law excluded from
$\params$, and $\params$-free likelihood factors --- the last licensing observations at several
noise levels to add into one $\Xis$. Given those, maximisation costs $O(P\ngrid^2)$ for $P$
parameters, independently of $\nseq$, while forming $\Xis$ costs $O(\nseq\nsites\ngrid^2)$ per
iteration. ``One matrix'' is therefore a statement about sufficiency, not about inference being
cheap.

The reported numbers come from one frozen configuration, defined in a single file and imported by
the experiments that produce them: Laplace chain unless stated, $\corr=0.85$, $\nsites=32$,
$\ngrid=401$, $\halfwidth=8$, twelve log-spaced noise levels on $[0.05,3.0]$, sixteen replicates,
and $256$ held-out sequences on a disjoint seed.

Two qualifications limit how much the frozen file guarantees on its own. A declared field is not
necessarily a \emph{read} field: importing a configuration does not make the code consult it, and
the provenance stamp compares the configuration against itself, so it cannot see a setting nobody
reads. Two such fields exist today and are listed as defects in
\texttt{tests/test\_frozen\_config\_is\_read.py}, which fails the build if a third appears.
Separately, experiments carrying their own local settings --- chain length, noise grid,
architecture sweeps --- are not described by this file at all, so their outputs additionally record
a digest of the settings actually resolved at run time.

\paragraph{What Table~\ref{tab:pointwise} covers.} Four qualifications, kept here rather than in
the caption. Both arms select their optimisation budget on a disjoint validation bundle --- the
network its parameterisation \emph{and} training length, EM--BP its iteration count --- and that
choice agrees with the test-set optimum in $\netagree\%$ and $\emagree\%$ of cells respectively,
so the protocol is symmetric and the selection is close to an oracle without being one. The
$\nseq=\extendedn$ row uses a raised budget, the caps having been set for
$\nseq\le\nseqcalib$; rerunning $\nseq=\nseqcalib$ at that budget on the same seeds moves its
ratio by $\budgetdelta\pm\budgetdeltase$, so the caps are selected without binding.
$\ndropped$ of $\ncellstotal$ cells are excluded by the resolution gate, because the narrowest
fitted mixture component is under two grid cells wide; including them changes no ratio by more
than $\dropshift$. And the network is beaten in all but $2$ of the reported cells, both at the
smallest sizes --- Appendix~\ref{app:aggregation} resolves where the advantage is largest.

\paragraph{An excluded row, and why the neighbouring one stays.} No $\nseq=8192$ row is reported,
although one was run. Its parameter file records a source tree with $207$ uncommitted paths, and
intersecting that list with the experiment's fifteen-file import closure leaves three that the run
itself executes: the experiment, the frozen configuration, and the EM implementation. The program
that produced those numbers is not recoverable, so the row is excluded from
Table~\ref{tab:pointwise} and from every figure drawn on the same grid; it would otherwise have
set the top of the quoted ratio range.

The $\nseq=4096$ rows are kept, and the distinction is not special pleading. Their tree carried
eight uncommitted paths and the same intersection is \emph{empty} --- all eight are wavelet and
image-metric work, none of it reachable from this experiment --- so no edit in that tree could have
moved these numbers. ``The tree was dirty'' and ``a file this run executes was dirty'' are
different facts, and only the second is disqualifying. The check is
\texttt{tools/provenance\_gate.py}, which computes the closure statically and reports the
intersection per run rather than refusing every dirty tree alike.

\paragraph{One view per chain.} Each clean chain is noised at \emph{one} level, the chains being
split evenly across the twelve levels, so a training set of $\nseq$ chains yields exactly $\nseq$
observations and the likelihood is a genuine one-view marginal likelihood rather than a composite
objective over repeated views of the same latent sequence. This costs resolution at the small end
--- at $\nseq=32$ a level carries two or three chains --- and it is what makes $\nseq$ a count of
independent latent sequences, so that the horizontal axis of Table~\ref{tab:pointwise} is a
sample size and not a count of noised views.

\paragraph{The stopping rule, and what actually gates a reported number.} Two mechanisms are in
play and only the second determines any headline result.

\texttt{fit\_em} stops when the signed log-evidence gain per transition falls in $[0,10^{-9}]$ ---
signed, so a decrease is recorded as a defect rather than accepted as convergence --- subject to an
iteration cap supplied by the caller, after which the run is labelled \emph{censored} rather than
converged. This is a diagnostic: it says whether a fit ran to a flat likelihood, and nothing more.

The headline estimator does not rely on it. It selects an iteration count on the disjoint
validation bundle, from the grid listed below, exactly as the network arm selects its training
length --- so the reported model is chosen by held-out risk, not by a convergence test, and a
censored fit is not thereby disqualified.

That distinction carries real weight here, because the coordinates converge at very different
rates: at $\corr=0.85$, $\nseq=512$, $C=8$ the autoregressive coefficient is flat from iteration
${\sim}25$ while the fitted excess kurtosis reads $0.84$ at $30$ iterations, $1.85$ at $60$,
$2.30$ at $120$ and $2.29$ at $400$. A likelihood- or $\corr$-based stopping rule would certify a
kernel still missing a third of the higher-order structure. Validation selection is what guards
against that, and it is the guard every reported number passes through. The configuration file
also declares a shape-based threshold, $|\Delta\kappa_4|<10^{-3}$; nothing reads it, it is one of
the two unread fields recorded above, and no reported number was produced by it.

\begin{center}
\small
\begin{tabular}{ll}
\toprule
chain & Laplace AR(1), $\corr=0.85$, $\ivar=1-\corr^2$, $a_1\sim\normal{0}{1}$, $\nsites=32$ \\
other families & Gaussian, Student-$t$ ($\nu=5$), uniform, symmetric two-component mixture
                 ($\kappa=0.6$) \\
                & all variance-matched to $\ivar$, so $\Cov(a_i,a_j)=\corr^{|i-j|}$ exactly \\
grid & $\ngrid=401$, $\halfwidth=8$ ($\gridh=0.04$); refinements $\ngrid\in\{801,1601\}$ \\
noise levels & twelve log-spaced $t\in[0.05,3.0]$, for training and pointwise evaluation \\
budgets & $\nseq\in\{32,64,128,256,512,1024,2048\}$, one noise view per chain \\
test set & $256$ held-out sequences, seed \texttt{exp07-test}, disjoint from all training seeds \\
validation & a second disjoint bundle, used to select both arms' budget (below) \\
replicates & sixteen; the replicate index is mixed into every training seed, so replicates \\
           & vary the training set \emph{and} the initialisation; the test set is common to all \\
\midrule
EM--BP & mixture-innovation kernel \eqref{eq:kernel}, $\ncomp=8$; $3\ncomp=24$ free
         parameters \\
       & $\corr$ initialised at $0.3$; iteration count selected on validation from \\
       & $\{10,20,40,60,80,120,160,220,300,400\}$; relative tolerance $10^{-9}$ \\
global MLP & fully connected, $25{,}248$ parameters, denoising score matching, Adam \\
           & training length selected on validation from
             $\{0.5,1,2,3.5,6,10,14,20\}\times10^{3}$ steps \\
           & both parameterisations ($\varepsilon$ and $a$) trained, and the choice between
             them \\
           & is also made on validation \\
local CNN & weight-shared one-dimensional convolution, radius $r$, receptive field $2r+1$ \\
capacity control & $24$ configurations, $3.2$k--$297$k parameters \\
\midrule
hardware & CPU: single-threaded, \texttt{medium\_cpu} partition, Bocconi HPC \\
         & GPU: CUDA 12.4, CuPy 13.6, gated by a CPU/GPU parity test before any sweep \\
\bottomrule
\end{tabular}
\end{center}

\paragraph{Selection rules, stated plainly.} Every tunable quantity that either arm is allowed to
choose is chosen on a validation bundle disjoint from both the training set and the test set, and
the two arms are given the same privilege. The network selects its parameterisation
($\varepsilon$ or $a$) \emph{and} its training length; EM--BP selects its iteration count. Neither
selection touches the test set. The selected choice coincides with the test-set optimum in $82.2\%$
of cells for the network and $79.5\%$ for EM--BP, and Table~\ref{tab:pointwise} reports the
validation-selected numbers, not those optima.

The symmetry matters because the training length is a bias--variance knob for both arms, not a
convergence detail: selecting one arm's budget while pinning the other's is not a neutral choice.
Measured on this run, giving the network the same early-stopping privilege costs between $5\%$ and
$36\%$ of the ratio, the larger corrections at the smaller sample sizes. The one selection that
remains a fixed configuration value rather than a validation choice is the convolution's receptive
field; the numbers in \S\ref{sec:networks} that depend on it are labelled exploratory.

\section{Scope conditions}
\label{app:limits}

The initial law is held fixed at the one the generating process uses, so nothing is misspecified
there and nothing shows it could be recovered. The linear form of \eqref{eq:kernel} is assumed:
what is free is a scalar and a density, not an arbitrary bivariate transition. The chain is
covariance-stationary rather than strictly stationary. Identifiability is a theorem at level 1 and
conditional at level 2. The implementation is a discretisation whose error, for the Laplace family,
is second order rather than spectral. Inference costs $O(\nsites\ngrid^{2})$ per sequence, far above
a network forward pass, so the estimator's advantage is in data and not in wall-clock. And the
comparison is asymmetric in supplied information by construction
(Section~\ref{sec:networks}): it measures what the Markov assumption buys where that assumption
holds, and says nothing about a setting where it does not.

\section{Reproducibility map}
\label{app:repro}

\begin{center}
\scriptsize
\setlength{\tabcolsep}{4pt}
\renewcommand{\arraystretch}{1.15}
\begin{tabular}{@{}>{\raggedright\arraybackslash}p{0.30\textwidth}
                   >{\raggedright\arraybackslash}p{0.30\textwidth}
                   >{\raggedright\arraybackslash}p{0.34\textwidth}@{}}
\toprule
claim & script (\texttt{experiments/\allowbreak }) & output (\texttt{outputs/\allowbreak }) \\
\midrule
grid and domain convergence & \texttt{exp\_\allowbreak 01\_\allowbreak grid\_\allowbreak validation} &
  \texttt{exp\_\allowbreak 01\_\allowbreak .../\allowbreak grid\_\allowbreak heatmap.csv} \\
boundary and quadrature diagnostics & \texttt{exp\_18\_... -{}-parts boundary} &
  \texttt{exp\_\allowbreak 18/\allowbreak boundary.csv} \\
closure error vs $t$ & \texttt{exp\_\allowbreak 02\_\allowbreak laplace\_\allowbreak gaussian\_\allowbreak ...} &
  \texttt{exp\_\allowbreak 02\_\allowbreak .../\allowbreak laplace\_\allowbreak summary.csv} \\
closure error by family & \texttt{exp\_\allowbreak 03\_\allowbreak nongaussian\_\allowbreak ...} &
  \texttt{exp\_\allowbreak 03\_\allowbreak .../\allowbreak innovation\_\allowbreak sweep.csv} \\
EM monotonicity, $\corr$ recovery & \texttt{exp\_18\_... -{}-parts emtrace} &
  \texttt{exp\_\allowbreak 18/\allowbreak em\_\allowbreak trace.csv} \\
learned innovation density & \texttt{exp\_18\_... -{}-parts density} &
  \texttt{exp\_\allowbreak 18/\allowbreak innovation\_\allowbreak density.csv} \\
parameter-recovery rate & \texttt{exp\_\allowbreak 06\_\allowbreak em\_\allowbreak parameter\_\allowbreak ...} &
  \texttt{exp\_\allowbreak 06\_\allowbreak .../\allowbreak em\_\allowbreak rate.csv} \\
Fisher information vs $t$ & \texttt{exp\_\allowbreak 06\_\allowbreak em\_\allowbreak parameter\_\allowbreak ...} &
  \texttt{exp\_\allowbreak 06\_\allowbreak .../\allowbreak price\_\allowbreak of\_\allowbreak noising.csv} \\
Fisher identity vs finite diff. & \texttt{exp\_\allowbreak 08\_\allowbreak gradient\_\allowbreak vs\_\allowbreak ...} &
  \texttt{exp\_\allowbreak 08\_\allowbreak .../\allowbreak gradient\_\allowbreak vs\_\allowbreak exact.csv} \\
sample efficiency (\S\ref{sec:networks}) & \texttt{exp\_\allowbreak 07\_\allowbreak em\_\allowbreak vs\_\allowbreak score\_\allowbreak ...} &
  \texttt{replicates/\allowbreak merged\_\allowbreak summary.csv} \\
network capacity control & \texttt{exp\_\allowbreak 07\_\allowbreak em\_\allowbreak vs\_\allowbreak score\_\allowbreak ...} &
  \texttt{exp\_\allowbreak 07\_\allowbreak .../\allowbreak capacity.csv} \\
inference and training cost & \texttt{exp\_\allowbreak 07\_\allowbreak em\_\allowbreak vs\_\allowbreak score\_\allowbreak ...} &
  \texttt{exp\_\allowbreak 07\_\allowbreak .../\allowbreak \{inference,training\}\_cost.csv} \\
receptive field (exploratory) & \texttt{exp\_\allowbreak 12\_\allowbreak receptive\_\allowbreak field} &
  \texttt{exp\_\allowbreak 12\_\allowbreak .../\allowbreak efficiency\_\allowbreak three\_\allowbreak way.csv} \\
sampler step convergence & \texttt{exp\_\allowbreak 16\_\allowbreak ... MODE=calibrate} &
  \texttt{exp\_\allowbreak 16/\allowbreak calibrate\_\allowbreak steps*/\allowbreak steps.csv} \\
generation by budget & \texttt{exp\_\allowbreak 16\_\allowbreak ... MODE=generate} &
  \texttt{exp\_\allowbreak 16/\allowbreak gen\_\allowbreak n*\_\allowbreak seed*/\allowbreak generation.csv} \\
generation by family & \texttt{exp\_\allowbreak 16\_\allowbreak sampling\_\allowbreak validation} &
  \texttt{exp\_\allowbreak 16/\allowbreak family\_\allowbreak */\allowbreak generation.csv} \\
capacity: generation & \texttt{exp\_\allowbreak 16\_\allowbreak sampling\_\allowbreak validation} &
  \texttt{exp\_\allowbreak 16/\allowbreak components\_\allowbreak C*/\allowbreak generation.csv} \\
capacity: pointwise & \texttt{exp\_\allowbreak 16\_\allowbreak sampling\_\allowbreak validation} &
  \texttt{exp\_\allowbreak 16/\allowbreak cpoint\_\allowbreak C*/\allowbreak pointwise.csv} \\
discrete-alphabet variant & \texttt{exp\_\allowbreak 10\_\allowbreak discrete\_\allowbreak alphabet} &
  \texttt{exp\_\allowbreak 10\_\allowbreak .../\allowbreak discrete\_\allowbreak *.csv} \\
\bottomrule
\end{tabular}
\end{center}

Every figure in this note is produced by \codefile{paper/figures/make\_figures.py}, which reads
the outputs above and fails loudly if one is missing; no number in any figure is entered by hand.
Exact commands, environment and HPC job configuration are in \codefile{REPRODUCIBILITY.md}, and
the claim-level audit behind this revision is \codefile{REVISION\_AUDIT.md}.

\section{The structured baseline}
\label{app:structured}

The headline comparison of Section~\ref{sec:networks} uses a fully connected network, which
carries no locality prior and no weight sharing. Separating that from the structural assumption
requires a weight-shared window predictor with radius and parameterisation selected on a
validation bundle disjoint from test, and EM--BP refitted on the same chains --- at the headline
protocol, with every site scored, deployed and provenance-gated through
\texttt{hpc/deploy\_clean.sh}. \texttt{experiments/exp\_31\_structured\_baseline.py} is that
measurement: it replaces an earlier exploratory run whose provenance did not verify and whose
protocol differed from the headline table's in four ways (\texttt{outputs/README\_exp12\_scaled.md}
records the detail; no number from that run is cited here or in the main text).

\paragraph{Which architecture, and under what search.} Three candidates were screened on
$\screendevseeds$ development seeds, disjoint from the $\structseeds$ used below: the
weight-shared window head, a dilated convolutional stack, and a bidirectional message-passing
network, the last two able in principle to propagate across the whole chain. Each ran over its own
hyperparameter grid, with a shared learning-rate grid, both output parameterisations and a common
checkpoint ladder ($\screenckptlo$ to $\screenckpthi$ steps), $\screenconfigs$ configurations in
all, selected on mean risk over development seeds on the \screenregion{} region. \textbf{The window
head won}, at radius $\screenwinnerradius$, width $\screenwinnerwidth$.

%% GENERATED by tools/make_tab_screening.py from outputs/frozen/exp_31_screen
%% (4 development seeds, 768 scored rows, provenance-clean).
%% Do not hand-edit the numbers; rerun the generator.

\begin{center}\small
\captionof{table}{The architecture screen that selected Table~\ref{tab:structured}'s
baseline. Each candidate architecture was run over its own hyperparameter grid on
4 development seeds (disjoint from the confirmatory seeds),
at $\nseq \in \{32, 2048\}$, with a shared learning-rate grid,
both output parameterisations, and a shared checkpoint ladder; the reported risk is
the mean over development seeds on the interior slice used for selection. $\star$
marks the winner carried into the confirmatory run. The dilated stack carries an
order of magnitude more parameters than the winner and finishes last, so the
outcome is not an artefact of one arm being given more capacity than another.}
\label{tab:screening}
\begin{tabular}{lccc}
\toprule
architecture & configurations & parameters & selection risk \\
\midrule
weight-shared window head\;$\star$ & 16 & 4.8k--21.4k & $0.137$ \\
dilated convolutional stack & 8 & 66.4k--264k & $0.205$ \\
bidirectional message passing & 8 & 9.0k--34.3k & $0.172$ \\
\bottomrule
\end{tabular}
\end{center}


This is a model-selection fact under this search space and protocol, not evidence that long-range
propagation is intrinsically unnecessary or that radius $\screenwinnerradius$ is optimal in the
population. It does rule out the cheapest objection: the dilated stack carries
$\screenconvparams$ parameters against the window head's $\screenwindowparams$ and placed last, so
the long-range arms did not lose for want of capacity. The confirmatory comparison below is
therefore against the strongest baseline \emph{this screen found}, not the strongest conceivable
one, and every claim that follows is scoped to that.

%% GENERATED by tools/make_tab_structured.py from outputs/frozen/exp_31_confirm_merged
%% (16 seeds, 4 sizes, 384 rows, provenance-clean).
%% Do not hand-edit the numbers; rerun the generator.

\begin{center}\small
\captionof{table}{EM--BP against the screened structure-aware baseline, at the
headline protocol throughout. The baseline is a weight-shared window head whose
architecture and hyperparameters were fixed by the screen of
Table~\ref{tab:screening} on development seeds; only the checkpoint is chosen
here, on a validation bundle disjoint from test, and EM--BP's iteration count is
chosen the same way. Risks are all-site and averaged over the 16 seeds; the ratio is
risk(window)/risk(EM--BP) formed within each seed and then averaged, $\pm$ one
s.e. \emph{All} scores every site, and is the headline; \emph{centre} is a
single-site diagnostic. A third region, the predeclared interior slice, is scored
in the stored output and used for architecture selection but not reported here.}
\label{tab:structured}
\begin{tabular}{rcccc}
\toprule
 & \multicolumn{2}{c}{risk, all sites} & \multicolumn{2}{c}{ratio} \\
$\nseq$ & EM--BP & window & all sites & centre site \\
\midrule
32   & $0.1005$ & $0.2006$ & $2.34 \pm 0.26$ & $2.04 \pm 0.24$ \\
128  & $0.0652$ & $0.1594$ & $2.82 \pm 0.30$ & $2.35 \pm 0.26$ \\
512  & $0.0394$ & $0.1353$ & $3.58 \pm 0.18$ & $3.25 \pm 0.18$ \\
2048 & $0.0224$ & $0.1267$ & $6.14 \pm 0.45$ & $5.43 \pm 0.39$ \\
\bottomrule
\end{tabular}
\end{center}


EM--BP is more accurate in every one of the $\structcells$ scored cells --- $\structseeds$ seeds
$\times$ $\structsizes$ sizes $\times$ $\structregions$ regions (\structregionlist{}), of which the
table shows two. That count is descriptive: regions and sizes share fitted models and data, so the
inferential unit behind each row is the $\structseeds$ paired seeds. The ratio increases across the
four sizes tested. That is a description of four measurements and not evidence that the residual is
non-architectural: a baseline with an irreducible approximation floor $b^2$ gives
$(b^2 + c_w/\nseq)/(c_e/\nseq)$, linear increasing in $\nseq$ for a wholly architectural reason,
and the raw risks above are consistent with it --- EM--BP's risk falls about fourfold across the
range while the window head's falls by less than half that. Both arms select their budget
symmetrically on a validation bundle disjoint from test, \emph{within the grids on offer}:
$\structcapem\%$ of EM--BP cells and $\structcapwindow\%$ of window-head cells take the last rung,
which indicates validation error still falling where the grid ended rather than convergence. The
values therefore compare the two procedures at stated, validation-selected budgets, not at
converged optima. Appendix~\ref{app:aggregation} reports the ratio under four aggregations of the
same cells; all exceed $2\times$ at every size.

\section{Density-level recovery: the Hellinger figures}
\label{app:hellinger}

Both kernels are put on the same grid and renormalised under the same quadrature, so the
comparison measures fit rather than normalisation convention. Absent a law over parent states a
conditional kernel $\kernel(\cdot\mid a)$ has no canonical scalar distance, and an unweighted
median over grid columns gives a rarely-visited tail parent the same standing as a typical one ---
so it was, until this revision, a diagnostic on the fitted density rather than an intrinsic
distance between transition laws. The intrinsic version is $\E_{A\sim\pi_0}[H(\kernel_{\hat\params}
(\cdot\mid A), \kernel_0(\cdot\mid A))]$ under the true parent law $\pi_0$, and both are reported
below, computed by the same run (\texttt{outputs/exp\_07\_certified\_clean\_seed*/}, sixteen
seeds, clean-deployed and gated; independently verified bit-identical to the certified table's
own score-based fields, \texttt{outputs/README\_exp07\_certified\_clean.md}).

%% GENERATED by tools/make_tab_hellinger.py -- do not hand-edit.
%% Source: outputs/exp_07_certified_clean_seed*/ (16 seeds, clean
%% deploy, verified bit-identical to outputs/frozen/exp_07_certified_seed*/
%% on every score-based field -- see outputs/README_exp07_certified_clean.md.

\begin{center}\small
\begin{tabular}{@{}lccccccc@{}}
\toprule
$\nseq$ & $32$ & $64$ & $128$ & $256$ & $512$ & $1024$ & $2048$ \\
\midrule
relative score error & $0.0580$ & $0.0501$ & $0.0342$ & $0.0258$ & $0.0220$ & $0.0162$ & $0.0132$ \\
transition Hellinger, median & $0.1030$ & $0.0896$ & $0.0757$ & $0.0637$ & $0.0598$ & $0.0522$ & $0.0500$ \\
transition Hellinger, weighted mean & $0.0994$ & $0.0867$ & $0.0741$ & $0.0630$ & $0.0592$ & $0.0521$ & $0.0499$ \\
\bottomrule
\end{tabular}
\end{center}


\noindent All three rows fall with $\nseq$, so the transition is being recovered and not merely
fitted. The weighted mean sits consistently a little below the unweighted median --- the parent
law downweights exactly the tail columns an unweighted vote overweights --- without changing the
qualitative story.

%% GENERATED by tools/make_aggregation_robustness.py -- do not hand-edit.

\section{Does the headline depend on how the ratio was averaged?}
\label{app:aggregation}

A ratio is a nonlinear function of two means, so its value depends on the order
of averaging: $\E[X/Y] \ge \E[X]/\E[Y]$ by Jensen, with a gap that grows
with the spread. Table~\ref{tab:pointwise} reports the mean of per-cell
ratios, aggregated per seed. That is the right inferential unit --- the twelve
noise levels share a training set and a fitted model --- but being the right
unit does not make it the only defensible summary, so here is the same data
under five.

\begin{center}\small
\captionof{table}{The headline under five estimands. (1) is what
Table~\ref{tab:pointwise} reports. The interval is a paired bootstrap over
seeds, 20,000 resamples; seeds are resampled rather than cells, because cells
within a seed share a fitted model.}
\label{tab:aggregation}
\begin{tabular}{rccccc}
\toprule
$\nseq$ & (1) per-cell & (2) ratio of means & (3) geometric & (4) median
        & (5) paired 95\% CI \\
\midrule
32 & $7.0$ & $6.8$ & $6.1$ & $6.8$ & $[6.1,\,8.1]$ \\
64 & $7.8$ & $6.6$ & $6.7$ & $6.6$ & $[6.3,\,9.5]$ \\
128 & $8.8$ & $8.0$ & $7.8$ & $8.6$ & $[7.9,\,9.6]$ \\
256 & $9.7$ & $8.8$ & $8.6$ & $9.2$ & $[8.5,\,10.9]$ \\
512 & $9.0$ & $8.6$ & $7.9$ & $8.2$ & $[8.1,\,9.9]$ \\
1024 & $12.5$ & $9.7$ & $10.8$ & $11.1$ & $[10.5,\,14.7]$ \\
2048 & $15.5$ & $10.0$ & $13.2$ & $12.9$ & $[13.0,\,18.9]$ \\
4096 & $17.6$ & $12.4$ & $15.6$ & $16.0$ & $[15.4,\,19.9]$ \\
8192 & $20.2$ & $16.1$ & $18.8$ & $19.5$ & $[18.3,\,22.0]$ \\
\bottomrule
\end{tabular}
\end{center}

\noindent Two things follow, and the second is the one worth stating plainly.

The conclusion does not depend on the choice: every estimand at every size
exceeds $6.1$, every bootstrap interval sits well clear of $1$, and no
aggregation produces a reversal anywhere.

But \textbf{the estimand we report is systematically the largest of the
four}, and by a margin that grows with $\nseq$ --- up to $5.5$ at
$\nseq=2048$. The ratio of per-seed averaged errors, which is the summary a
reader is most likely to have in mind, runs $6.6$ to $16.1$
rather than $\ratiolo$ to $\ratiohi$. This is Jensen's inequality doing
exactly what it must, not an error, but quoting only the largest of five
summaries would be a choice presented as a fact.

Resolving the average over noise levels shows where the aggregation hides
structure: the ratio is not flat in $t$ but peaks near $t \approx 0.5$ and
falls at both ends, and the weakest single (size, level) cell in the whole
grid is $3.6$. The estimator's advantage is largest at moderate
noise, where the posterior is neither dominated by the likelihood nor by the
prior --- which is where knowing the transition structure should help most.


The same question applies to the structure-aware ratio, and with more force, since that one is
quoted as a range. There the noise schedule is already reduced into each cell's risk, so the seed
is the inferential unit by construction and only the forming and pooling of the ratio remain
optional.

%% GENERATED by tools/make_aggregation_structured.py from outputs/frozen/exp_31_confirm_merged
%% (region=all, 16 seeds, 4 sizes, provenance-clean).
%% Do not hand-edit the numbers; rerun the generator.

\begin{center}\small
\captionof{table}{Table~\ref{tab:structured}'s all-site ratio under four
aggregations of the same cells, with the paired bootstrap interval over seeds.
(1) is what the table reports. The conclusion does not depend on the choice:
every estimand at every size exceeds $1.98\times$, and the largest
departure from the reported value is $0.48$.}
\label{tab:aggregation-structured}
\begin{tabular}{rccccc}
\toprule
$\nseq$ & (1) per-seed & (2) ratio of means & (3) geometric & (4) median & (5) 95\% CI on (1) \\
\midrule
$32$ & $2.34$ & $2.00$ & $2.15$ & $1.98$ & $[1.88,\, 2.86]$ \\
$128$ & $2.82$ & $2.45$ & $2.61$ & $2.62$ & $[2.29,\, 3.44]$ \\
$512$ & $3.58$ & $3.43$ & $3.51$ & $3.59$ & $[3.25,\, 3.91]$ \\
$2048$ & $6.14$ & $5.66$ & $5.89$ & $5.93$ & $[5.30,\, 7.02]$ \\
\bottomrule
\end{tabular}
\end{center}

